\begin{table}[H]
\centering
\caption{Operational thresholds and their roles. The regional-mean monthly accumulation is the headline FLOOD definition; the mine-cell daily variant is a conservative point-level lower bound. Source anchors are listed in the footnote.}
\label{tab:thresholds}
\small
\begin{tabular}{p{3.0cm}p{3.8cm}p{2.8cm}p{4.4cm}}
\toprule
\textbf{Channel} & \textbf{Threshold} & \textbf{Resolution} & \textbf{Role} \\
\midrule
\texttt{FLOOD\_regmean\_M}$^{1}$ & >= 20 mm/month & monthly & FLOOD headline \\
\texttt{FLOOD\_mine}$^{2}$ & >= 20 mm/day & daily & Conservative point-level lower bound \\
\texttt{HEAT\_mine}$^{3}$ & >= 23.5 degC monthly mean & monthly aggregate & Acute (robust null) \\
\texttt{DROUGHT\_meteo}$^{4}$ & SPEI-12 <= -1 & monthly aggregated to annual share & DROUGHT headline (meteorological) \\
\bottomrule
\end{tabular}
\\[2pt]
{\footnotesize\raggedright
$^{1}$\,Wilcox et al. 2016 GRL 20 mm threshold applied at the regional-mean monthly scale.\quad
$^{2}$\,Wilcox et al. 2016 GRL aluvion definition.\quad
$^{3}$\,Chile DS 594/1999 occupational heat-stress (daily 27 C anchor calibrated to monthly mean on CR2MET 1960-2021).\quad
$^{4}$\,WMO moderate-drought standard (Vicente-Serrano 2010).
}
\end{table}

